\section{Proposed Experimental Validation}
\label{sec:experiments}

\subsection{Datasets}
\label{subsec:datasets}

The validation framework targets four distinct benchmark datasets: TartanAir~\cite{wang2020tartanair}, Blackbird, MARSIM~\cite{kong2023marsim}, and AeroVerse. Each dataset is converted to the common \texttt{SequenceData} contract with immutable source version, license record, coordinate convention, calibration metadata, and SHA-256 manifest. The cohort identifier \texttt{frozen-2026-01} denotes the study cohort rather than a vendor release and resolves to exact source versions and file hashes.

Data are split at the environment or flight level, never across adjacent temporal frames. The training split is the only source of learned weight updates; calibration splits are strictly reserved for hyperparameter thresholds, normalization, early stopping, and uncertainty calibration; the test split remains inaccessible until code, artifacts, and analysis scripts are frozen.

The preregistered offline mechanism study uses 20 preselected test sequences per dataset, yielding 80 paired sequence-level units. Systems validation additionally defines 30 paired Software-in-the-Loop (SIL) scenarios with five random seeds, 20 Hardware-in-the-Loop (HIL) scenarios with five seeds, and 20 paired controlled real-flight missions per variant when safety gates remain open.

A dedicated UAV dataset remains desirable for synchronized RGB, thermal, IMU, GNSS, point cloud, elevation/reference maps, static/dynamic ground truth labels, visibility masks, and controlled sensor degradation scenarios. Metrics that require unavailable ground truth are reported as unavailable rather than silently omitted or imputed.

\subsection{Two-Level Comparison Design}
\label{subsec:comparison_design}

The evaluation is decoupled into two non-interchangeable levels. Both levels use common dataset and metric contracts where applicable, but they address distinct scientific questions and produce separate statistical reports.

\subsubsection{External System Comparison}
\label{subsubsec:external_comparison}

The external comparison evaluates whether the complete S4D-TAM system is competitive with independently implemented navigation and SLAM systems under identical data, hardware constraints, sensor availability, and evaluator definitions.

The mandatory primary baseline set comprises:
\begin{itemize}
  \item \textbf{ORB-SLAM3},
  \item \textbf{VINS-Mono}~\cite{qin2018vinsmono},
  \item \textbf{FAST-LIO2}~\cite{xu2022fastlio2},
  \item \textbf{LIO-SAM}~\cite{shan2020liosam}.
\end{itemize}
Each baseline is pinned to an upstream git commit revision, fixed containerized environment, calibration profile, sensor-topic mapping, loop-closure policy, warm-up policy, and hardware execution budget. Results enter the benchmark exclusively through the common \texttt{AlgorithmResult} schema.

This level measures complete-system competitiveness. Observed differences relative to external baselines cannot be interpreted as the isolated causal effect of an individual S4D-TAM component.

\subsubsection{Internal Mechanism Study}
\label{subsubsec:internal_study}

The internal comparison isolates mechanism contributions using the full system and seven single-component ablation variants:
\begin{itemize}
  \item \texttt{H1\_no\_semantics},
  \item \texttt{H2\_no\_temporal\_state},
  \item \texttt{H3\_no\_calibrated\_uncertainty},
  \item \texttt{H4\_no\_topology},
  \item \texttt{H5\_no\_reference\_map},
  \item \texttt{H6\_no\_risk\_prediction},
  \item \texttt{H7\_no\_token\_lifecycle}.
\end{itemize}
Each variant differs from the full system by exactly one switch and utilizes its own frozen trained artifact. External systems are excluded from this matrix to avoid confounding complete-system identity with the causal effect of removing an individual mechanism.

\subsection{Metrics}
\label{subsec:metrics}

Primary endpoints are preregistered:
\begin{itemize}
  \item \textbf{Localization}: Absolute Trajectory Error (ATE) RMSE after $\mathrm{SE}(3)$ rigid alignment without scale optimization.
  \item \textbf{Forecasting}: Occupancy IoU at $1\,\text{s}$ (and $3\,\text{s}$ for multi-step characterization).
  \item \textbf{Safety}: Mission success rate and collision count per kilometer.
  \item \textbf{Efficiency}: 95th percentile (p95) processing latency and persistent map-memory footprint.
\end{itemize}

Secondary localization metrics include Relative Pose Error (RPE) translation RMSE, final drift in meters and percentage of traveled distance, rotation RPE, relocalization success rate, and time-to-relocalize. Semantic metrics encompass mIoU and macro F1 score. Forecasting secondary metrics include flow End-Point Error (EPE), Brier score, Negative Log-Likelihood (NLL), and Expected Calibration Error (ECE). Navigation metrics include near misses, minimum obstacle clearance, and path efficiency.

\subsection{Confirmatory Protocol and Statistical Analysis}
\label{subsec:statistical_analysis}

The preregistered H1--H7 mechanism study conducts sequence- or mission-level inference rather than treating individual frames as independent samples. Continuous outcomes are evaluated via paired mean differences and $95\%$ BCa bootstrap confidence intervals using $10{,}000$ resamples. Mission success is modeled as a paired scenario-level binary probability outcome, while distance-normalized event counts utilize Poisson or negative-binomial generalized linear models with log-distance offset.

The seven H1--H7 confirmatory decisions form one family controlled by the sequential Holm step-down procedure at family-wise error rate $\alpha = 0.05$. Hypotheses H1--H6 require the adjusted $p$-value and confidence interval to demonstrate superiority in the preregistered direction. Hypothesis H7 additionally requires non-inferiority in mission success within a $2$ percentage-point margin while demonstrating p95 latency reduction.
